\section*{Chapitre \ref{ch:b2:metamorphic-facies} --- Réactions métamorphiques et faciès}

\begin{solution}{exo:b2:metamorphic-facies:1}
(a) \omterm{def:b2:metamorphic-facies:metamorphism}{Métamorphisme} de contact, \omterm{def:b2:metamorphic-facies:metamorphism}{métamorphisme} thermique local dû à la chaleur de l'intrusion.
(b) \omterm{def:b2:metamorphic-facies:metamorphism}{Métamorphisme} de dislocation, dans une zone de faille~: réduction de la taille des grains, une mylonite.
(c) \omterm{def:b2:metamorphic-facies:metamorphism}{Métamorphisme} océanique, dû à l'eau de mer chaude près de la dorsale, avec une
recristallisation incomplète. (d) \omterm{def:b2:metamorphic-facies:metamorphism}{Métamorphisme} d'enfouissement, régional, sans déformation ni
schistosité.
\end{solution}

\begin{solution}{exo:b2:metamorphic-facies:2}
$\Delta_r S = 95.40 - 82.80 = \qty{12.60}{J.K^{-1}}$, $\Delta_r V = 4.986 - 4.415 =
\qty{0.571}{J.bar^{-1}}$~; pente $12.60/0.571 = \qty{22.1}{bar/K} = \qty{2.21}{MPa/K}$.
Sur \qty{100}{\celsius}, la pression d'équilibre augmente de \qty{2.21}{kbar}.
\end{solution}

\begin{solution}{exo:b2:metamorphic-facies:3}
(a) Faciès des schistes verts. (b) Faciès des schistes bleus. (c) Amphibolite~: hornblende avec plagioclase à plus de
\qty{17}{\percent} d'anorthite. (d) Éclogite.
\end{solution}

\begin{solution}{exo:b2:metamorphic-facies:4}
(a) $\rho = M/V$~: disthène $162.046/44.15 = \qty{3.670}{g.cm^{-3}}$, andalousite
$162.046/51.52 = 3.145$, sillimanite $162.046/49.86 = 3.250$. (b) Disthène~: à plus haute
pression, le terme $V\,(P - 1)$ de l'énergie de Gibbs croît le moins pour le plus petit volume.
(c) Les cristaux naturels contiennent d'autres éléments (le fer ferrique, par exemple, remplace une partie de
l'aluminium dans l'andalousite et la sillimanite naturelles), et les mesures sur des spécimens naturels sont
moins précises que les volumes molaires des tables.
\end{solution}

\begin{solution}{exo:b2:metamorphic-facies:5}
(a) $\Delta_r G = 3900 - 298.15\times8.59 = \qty{+1339}{J} > 0$~: disthène. (b) À
\qty{1}{bar}, $\Delta_r H = T\Delta_r S$~: $T = 3900/8.59 = \qty{454}{K} =
\qty{181}{\celsius}$. (c) $P = (8.59\times673.15 - 3900)/0.737 = 1882/0.737 =
\qty{2554}{bar} = \qty{2.55}{kbar}$, à comparer aux 2.54 mesurés. (d) La réaction est trop lente
à basse température pour se produire~: l'andalousite subsiste de façon métastable, comme la plupart des
minéraux de haute température à la surface.
\end{solution}

\begin{solution}{exo:b2:metamorphic-facies:6}
(a) $\Delta_r H = 7700 + 1000\times(121.39 - 122.29) = \qty{6800}{J}$~; $\Delta_r S = 296.44
- 285.02 = \qty{11.42}{J.K^{-1}}$~; $\Delta_r G = 6800 - 1000\times11.42 = \qty{-4620}{J}$.
(b) Avec les valeurs à \qty{298.15}{K}~: $7700 - 1000\times12.60 = \qty{-4900}{J}$.
(c) La sillimanite. La méthode simple se trompe de \qty{280}{J}, soit environ \qty{6}{\percent}~; rapportée
à $\Delta_r V = 0.571$, cela représente environ \qty{490}{bar} sur la droite, bien moins que l'effet de
l'incertitude des enthalpies.
\end{solution}

\begin{solution}{exo:b2:metamorphic-facies:7}
(a) $\Delta_r G = -1127.4 + 1128.5 = \qty{1.1}{kJ}$ et $\Delta_r V = (34.15 -
36.93)/10 = \qty{-0.278}{J.bar^{-1}}$~: $P = 1100/0.278 = \qty{3957}{bar} \approx
\qty{4.0}{kbar}$. (b) $\Delta_r S = 88.0 - 91.7 = -3.7$~; pente $-3.7/-0.278 =
\qty{13.3}{bar/K}$. (c) Avec $\Delta_r H = 0$, $P = T\Delta_r S/\Delta_r V = 13.31\,T$~:
à \qty{573.15}{K}, \qty{7629}{bar} $= \qty{7.63}{kbar}$. (d) L'aragonite des coquilles est
métastable à la surface~: les organismes la fabriquent, et elle persiste parce qu'elle ne se transforme
en calcite que lentement.
\end{solution}

\begin{solution}{exo:b2:metamorphic-facies:8}
(a) $-32.4/-1.698 = \qty{19.1}{bar/K}$. (b) $P = (T\Delta_r S - \Delta_r H)/\Delta_r V =
(-32.4\times773.15 + 5000)/(-1.698) = (-25050 + 5000)/(-1.698) = \qty{11808}{bar}
\approx \qty{11.8}{kbar}$. (c) $\sqrt{2.6^2 + 3.6^2 + 1.0^2} = \qty{4.55}{kJ}$ sur
$\Delta_r H$, soit $4550/1.698 = \qty{2.7}{kbar}$ en pression. (d) Aux pressions du
faciès éclogite, la composante albite du plagioclase est instable par rapport à
la jadéite et le quartz~; le sodium passe dans le pyroxène (omphacite), et le plagioclase
disparaît.
\end{solution}

\begin{solution}{exo:b2:metamorphic-facies:9}
(a) De 1 à \qty{2}{kbar}~: $1000/34 = \qty{29.4}{bar/K}$~; de 2 à 3~: $1000/25 =
\qty{40.0}{bar/K}$. (b) $\Delta_r S$ étant fixé, $\Delta_r V$ est inversement proportionnel à
la pente~: il a diminué d'un facteur $40.0/29.4 = 1.36$. (c) L'eau libérée~: son volume molaire
diminue quand la pression augmente, alors que les volumes des solides varient à peine.
\end{solution}

\begin{solution}{exo:b2:metamorphic-facies:10}
(a) À \qty{873.15}{K}~: disthène $=$ andalousite $(8.59\times873.15 - 3900)/0.737 =
\qty{4885}{bar}$~; disthène $=$ sillimanite $(12.60\times873.15 - 7700)/0.571 =
\qty{5782}{bar}$~; andalousite $=$ sillimanite $(4.01\times873.15 - 3800)/(-0.166) =
\qty{1799}{bar}$. (b) À \qty{4885}{bar}, sous la droite disthène $=$ sillimanite (5782),
la sillimanite est plus stable que le disthène, et au-dessus de la droite andalousite $=$ sillimanite
(1799), plus stable que l'andalousite~: le point se situe dans le champ de la sillimanite, et
l'équilibre disthène $=$ andalousite y est métastable. (c) À \qty{4}{kbar}~: entre 1799
et \qty{5782}{bar}, sillimanite. À \qty{7}{kbar}~: au-dessus de \qty{5782}{bar}, le disthène l'emporte sur la
sillimanite, et au-dessus de \qty{4885}{bar} il l'emporte sur l'andalousite~: disthène.
\end{solution}

\begin{solution}{exo:b2:metamorphic-facies:11}
(a) $z = P/(\rho g) = 3.75\times10^8/(2700\times9.81) = \qty{14.2}{km}$, incertaine de
$0.25\times10^8/26487 = \qty{0.94}{km}$. (b) $(504 - 10)/14.16 =
\qty{35}{\celsius.km^{-1}}$. (c) Les trois minéraux ne coexistent à l'équilibre qu'au
\omterm{def:b2:metamorphic-facies:invariant}{point invariant}~: le chemin de la roche doit y passer, et les réactions doivent
s'arrêter là. La plupart des roches passent à côté et gardent un ou deux polymorphes.
\end{solution}

\begin{solution}{exo:b2:metamorphic-facies:12}
(a) X~: $T = 10 + 10\times30 = \qty{310}{\celsius}$ (\qty{583.15}{K}); $P = 3000\times
9.81\times30000 = \qty{8.83e8}{Pa} = \qty{8.83}{kbar}$. Y~: $T = 10 + 30\times20 =
\qty{610}{\celsius}$ (\qty{883.15}{K}); $P = 2800\times9.81\times20000 = \qty{5.49}{kbar}$.
(b) X~: calcite $=$ aragonite à $13.31\times583.15 = \qty{7.76}{kbar} < 8.83$, aragonite~;
albite $=$ jadéite + quartz à $(32.4\times583.15 - 5000)/1.698 = \qty{8.18}{kbar} <
8.83$, jadéite + quartz~; disthène $=$ andalousite à $(8.59\times583.15 - 3900)/0.737 =
\qty{1.51}{kbar}$, en dessous de la température du point triple~: disthène. Y~: calcite $=$
aragonite à \qty{11.75}{kbar} $> 5.49$, calcite~; albite $=$ jadéite + quartz à
\qty{13.9}{kbar}, albite~; au-dessus du point triple, disthène $=$ sillimanite à
\qty{6.00}{kbar} $> 5.49$ et andalousite $=$ sillimanite à \qty{1.56}{kbar} $< 5.49$~:
sillimanite. (c) X, à \qty{10}{\celsius\per\km}, est le trajet le plus froid~: la jadéite et
l'aragonite, à seulement \qty{310}{\celsius}, le placent dans le faciès des schistes bleus. Y, trois fois
plus chaud, finit en sillimanite à \qty{610}{\celsius} et pression modérée, dans le faciès
amphibolite.
\end{solution}

\begin{solution}{pb:b2:metamorphic-facies:1}
\noindent\textbf{Partie I.}

\textbf{1.} $\Delta_r H = -2589.9 + 2593.8 = \qty{3.9}{kJ}$~; $\Delta_r S = 91.39 - 82.80 =
\qty{8.59}{J.K^{-1}}$~; $\Delta_r V = 5.152 - 4.415 = \qty{0.737}{J.bar^{-1}}$.

\textbf{2.} $\qty{7.7}{kJ}$~; $\qty{12.60}{J.K^{-1}}$~; $\qty{0.571}{J.bar^{-1}}$.

\textbf{3.} $\qty{3.8}{kJ}$~; $\qty{4.01}{J.K^{-1}}$~; $\qty{-0.166}{J.bar^{-1}}$~; chacune est
la seconde moins la première, puisque l'andalousite s'annule.

\textbf{4.} La sillimanite a la plus grande entropie (95.40)~: elle se situe du côté haute température
de ses deux droites. Le disthène a le plus petit volume (4.415)~: il se situe du côté haute
pression de ses deux droites. L'andalousite, au plus grand volume, occupe le côté
basse pression.

\textbf{5.} $8.59/0.737 = \qty{11.66}{bar/K}$~; $12.60/0.571 = \qty{22.07}{bar/K}$~;
$4.01/(-0.166) = \qty{-24.16}{bar/K}$.

\textbf{6.} La sillimanite a plus d'entropie mais moins de volume que l'andalousite~: $\Delta_r S > 0$
et $\Delta_r V < 0$. Élever la température favorise la sillimanite, élever la pression aussi,
si bien que la droite descend en pression quand la température monte.

\noindent\textbf{Partie II.}

\textbf{7.} $3900 - 298.15\times8.59 = \qty{+1339}{J}$~; $7700 - 298.15\times12.60 =
\qty{+3943}{J}$. Toutes deux positives~: le disthène est stable à la surface.

\textbf{8.} $T = \Delta_r H/\Delta_r S$~: $3900/8.59 = \qty{454.0}{K} = \qty{181}{\celsius}$~;
$7700/12.60 = \qty{611.1}{K} = \qty{338}{\celsius}$~; $3800/4.01 = \qty{947.6}{K} =
\qty{674}{\celsius}$.

\textbf{9.} En chauffant à \qty{1}{bar}, le disthène cède la place à l'andalousite à
\qty{181}{\celsius}~; à \qty{338}{\celsius} le disthène n'est plus stable, donc l'équilibre disthène $=$
sillimanite est métastable là~; l'andalousite cède la place à la sillimanite à
\qty{674}{\celsius}. Le premier et le troisième sont des équilibres stables.

\textbf{10.} $(8.59\times673.15 - 3900)/0.737 = \qty{2554}{bar} = \qty{2.55}{kbar}$,
à \qty{0.01}{kbar} de l'expérience.

\textbf{11.} $(12.60\times973.15 - 7700)/0.571 = 4562/0.571 = \qty{7989}{bar} =
\qty{7.99}{kbar}$, \qty{0.13}{kbar} au-dessus de l'expérience.

\textbf{12.} Le disthène $=$ andalousite monte de \qty{181}{\celsius} à \qty{1}{bar} avec la
pente \qty{11.7}{bar/K}~; l'andalousite $=$ sillimanite descend d'environ \qty{11.5}{kbar} à \qty{200}{\celsius}
jusqu'à \qty{1}{bar} à \qty{674}{\celsius}~; le disthène $=$ sillimanite, le plus raide des deux équilibres à pente positive, croise les deux
près de \qty{514}{\celsius} et \qty{3.9}{kbar}, chaque droite n'étant stable que d'un côté
du croisement seulement (le diagramme du chapitre).

\noindent\textbf{Partie III.}

\textbf{13.} $a_1 = 8.59/0.737 = \qty{11.655}{bar/K}$, $b_1 = 3900/0.737 =
\qty{5291.7}{bar}$~; $a_2 = 12.60/0.571 = \qty{22.067}{bar/K}$, $b_2 = 7700/0.571 =
\qty{13485.1}{bar}$.

\textbf{14.} $T = (b_2 - b_1)/(a_2 - a_1) = 8193.4/10.412 = \qty{786.92}{K} =
\qty{513.8}{\celsius}$.

\textbf{15.} $P = 11.655\times786.92 - 5291.7 = \qty{3880}{bar} = \qty{3.88}{kbar}$.

\textbf{16.} $(4.01\times786.92 - 3800)/(-0.166) = (3155.5 - 3800)/(-0.166) =
\qty{3882}{bar}$~: la même pression à l'arrondi des données près, comme l'exige la proposition du chapitre.

\textbf{17.} \qty{514}{\celsius} se situe à \qty{10}{\celsius} de \qty{504}{\celsius}, dans
l'incertitude de \qty{20}{\celsius}~; \qty{3.88}{kbar} se situe à \qty{0.13}{kbar} de 3.75,
dans 0.25. L'accord est bon.

\textbf{18.} Non~: \qty{536}{\celsius} se situe à \qty{32}{\celsius} de l'expérience, hors de son
incertitude. Négliger les capacités thermiques n'était pas l'erreur principale~: ce sont les enthalpies.

\noindent\textbf{Partie IV.}

\textbf{19.} $\Delta_r H$ du passage disthène $\to$ andalousite augmente de \qty{1000}{J}~; à
$T$ fixé, $P = (T\Delta_r S - \Delta_r H)/\Delta_r V$ diminue de $1000/0.737 = \qty{1357}{bar}$.

\textbf{20.} $\Delta T = -1357/(a_2 - a_1) = -1357/10.41 = \qty{-130}{K}$~: le point triple
se déplace vers \qty{384}{\celsius}.

\textbf{21.} La droite disthène $=$ sillimanite descend de $1000/0.571 = \qty{1751}{bar}$, et
le point triple se déplace de $+1751/10.41 = \qty{+168}{K}$.

\textbf{22.} Chaque enthalpie prise seule, variée de $\pm\qty{2.0}{kJ}$, déplace le point triple de
$\pm\qty{260}{K}$ (andalousite) ou $\pm\qty{336}{K}$ (sillimanite)~: la calorimétrie seule permet des
points triples sur plusieurs centaines de degrés.

\textbf{23.} Les enthalpies tabulées ne sont pas purement calorimétriques~: dans leurs
incertitudes, elles sont ajustées pour que les droites calculées s'accordent aux expériences réversées.
Les expériences fixent le diagramme~; les tables le portent, de façon cohérente, vers d'autres réactions.

\textbf{24.} Le point triple calculé à partir des tables à propriétés constantes se situe à
\textbf{\qty[text-series-to-math=true]{514}{\celsius}}.
\end{solution}
